\newpage
\section{Auxiliary Functions for Synthesis}\label{sec:tech:algo}

This section describes three auxiliary functions used for generator
synthesis. The first of these, $\normPlan$, converts a SRE formula
into a set of abstract traces. The second, $\termDerive$, generates test generators from a set of abstract traces. The third,
$\S{termTrace}$, generates a straightline program from a single refined abstract trace.

\begin{algorithm}[t!]
  \Procedure{$\normPlan(A) := $}{
  $A \leftarrow \Code{Minimize}(A)$\;
  \Match{$A$}{
    \lCase{$\emptyset$}{ \Return{$\emptyset$} }
    \lCase{$\epsilon$}{ \Return{$\{\epsilon\}$} }
    \lCase{$\msg{op}{\phi}$}{ \Return{$\{\msg{op}{\phi}^\bot\}$} }
    \lCase{$A_1 \seqA A_2$}{ \Return{$\{ \pi_1 \seqA \pi_2 ~|~ \pi_1 \in \normPlan(A_1) \land \pi_2 \in \normPlan(A_2) \}$}}
    \lCase{$A_1 \lorA A_2$}{ \Return{$\normPlan(A_1) \cup \normPlan(A_2)$}}
    \Case{$A^*$}{
      \Match{$\normPlan(A)$}{
        \lCase{$\emptyset$}{ \Return{$\emptyset$} }
        \lCase{$\bigcup \overline{\msg{op}{\phi}}$}{ \Return{$\{[\overline{\msg{op}{\phi}^\bot}]^*\}$} }
        \lCase{$\Pi$}{ \Return{$\{ \pi_1 \listconcat \pi_2 \listconcat \cdots \listconcat \pi_n ~|~ \pi_i \in \Pi \}$} }
      }
    }
  }
  }
  \caption{Abstract Trace Normalization}
  \label{algo:norm}
\end{algorithm}

\paragraph{Normalization}
The SRE normalization function $\normPlan{}$, shown in
\autoref{algo:norm}, first perform a standard symbolic automaton minimization step to remove intersections and complements, and then recursively translates an input SRE into a set of
abstract traces. It divides $A \lorA A'$ into multiple abstract
traces, and leaves the Kleene star over multiple singleton events
unchanged.

\begin{lemma}\label{lemma:norm-sound}[Normalization is sound]
  The normalized result has the same denotation as the input SRE, that
  is, for all SRE $A$ and set of traces $\{\overline{\pi}\}$, {\small
    $\denot{A} = \bigcup_{i} \denot{\pi_i}$}
\end{lemma}

\begin{algorithm}[t!]
    \Procedure{$\S{TermDerive}(C) := $}{
      $E \leftarrow \emptyset$\;
      \ForEach{$(\Gamma,\pi) \in C$} {
        $e_1 \leftarrow \S{DeriveTrace}(\Gamma, \pi)$\;
        $e_2 \leftarrow \S{SynRecursion}(\Gamma, \pi, \Code{UnrollingBound})$\;
        $E \leftarrow E \cup \{e_1, e_2\}$\;
      }
      \Return{$\bigoplus_{e \in E} e$}\;
    }
    \Procedure{$\S{DeriveTrace}(\Gamma, \pi) := $}{
      \Match{$\Gamma$}{
      \lCase{$[]$}{ \Return{$\S{DeriveTraceAux}(\pi)$} }
      \Case{$x{:}\urt{b}{\phi} \cons \Gamma'$}{
      \Return{$\zassume{\overline{x}}{\phi[\vnu \mapsto x]}{\S{DeriveTrace}(\Gamma', \pi)}$}\;
      }
    }
    }
    \Procedure{$\S{DeriveTraceAux}(\Pi) := $}{
      \Match{$\Pi$}{
        \lCase{$\epsilon$}{ \Return{$()$} }
        \lCase{$\Pi^* \seqA \pi'$}{ \Return{$\S{DeriveTraceAux}(\pi')$} }
        \Case{$\msgBGray{op}{\overline{x}}{\phi} \seqA \pi'$}{
          $\overline{x'} \leftarrow \Code{GetFreshNames}(\overline{x})$\;
          $\phi_1 \leftarrow \phi\msubst{x}{x'}$\;
          $e \leftarrow \S{DeriveTraceAux}(\pi')$\;
          \Return{$\zassume{\overline{x'}}{\phi_1}{e}$}
        }
        \Case{$\msgB{op}{\overline{x}\;y}{\phi} \seqA \pi'$}{
          $\overline{x'} \leftarrow \Code{GetFreshNames}(\overline{x})$\;
          $y' \leftarrow \Code{GetFreshName}(y)$\;
          $\phi_1 \leftarrow \exists y.\phi\msubst{x}{x'}$\;
          $\phi_2 \leftarrow \phi\msubst{x}{x'}[y\mapsto y']$\;
          $e \leftarrow \S{DeriveTraceAux}(\pi')$\;
          \Return{$\zassume{\overline{x'}}{\phi_1}{\zlet{y'}{\eff{op}\;\overline{x'}}{\sassert{\phi_2}; e}}$}
        }
      }
      }
    \caption{Term Derivation}
    \label{algo:term-derive}
  \end{algorithm}

  \paragraph{Term Derivation} The term derivation function
  \(\S{TermDerive}\) is shown in \autoref{algo:term-derive}. It first
  converts the input type context into \(\mykw{assume}\) statements
  that use the type qualifiers attached to each variable (line 5), and
  then derives the abstract trace with the help of the
  \(\S{DeriveTrace}\), which recursively transforms the abstract trace
  into a straightline generator program.  Any Kleene stars are dropped
  (line 9), ghost events are converted into \(\mykw{assume}\)
  expressions (line 14), and concrete events are converted into
  effectful operations (line 16-21).  Note that the arguments of the
  effectful operation are provided by $\mykw{assume}$ and the values
  they return are always checked by $\mykw{assert}$ (line 21).

\begin{algorithm}[t!]
  \Procedure{$\S{SynRecursion}(\Gamma, \Pi, \Code{UnrollingBound}) := $}{
     $(T, \{\pi_i\}) \leftarrow \S{MatchTemplate}(\Pi)$\;
     \If{$\forall i. 0 \leq i < \Code{UnrollingBound}. \exists \tau. \Gamma \vdash \Code{Unroll}(T, \{\pi_i\}, i) : \tau$}{
        $e_i \leftarrow \S{DeriveTrace}(\pi_i)$\;
        \Return{$T(\overline{e_i})$}
     }
  }
  \caption{Synthesis of Recursive Programs}
  \label{algo:syn-recursion}
\end{algorithm}

\paragraph{Recursion Synthesis} Kleene stars ($\Pi^*$) in abstract
traces are interpreted as recursion points. The synthesis algorithm
can synthesize a recursive program by expanding the Kleene star using
recursion. As shown in \autoref{algo:syn-recursion}, we employ a
template-based strategy, in which the algorithm matches the given
abstract trace with a template $T$, together with a set of abstract
traces $\{\pi_i\}$ that are matched from the orginal trace. As
dictated by \textsc{TFixInd}, the algorithm unrolls the recursive body
up to a bounded depth and invokes refinement loop (lines 3 in
\autoref{algo:syn-recursion}) to ensure that the unrolled abstract
traces are well-typed.  Finally, the synthesized program is
instantiated by the template with derived programs from these abstract
traces.

\begin{example}[Recursion Template]
  A representative template is a non-deterministic loop, whose form is
  $T \doteq \zlet{f}{\mykw{fix}\ f.\lambda ().\ () \oplus \boxed{e_1};
    f\;(); \boxed{e_2}}{\boxed{e_3};f\;();\boxed{e_4}}$, allowing
  $\Pi^*$ in the abstract trace to be matched. For instance, consider
  the abstract trace $\pi_\Code{stack}$ shown in
  Example~\ref{ex:refinement-process}: %
  \vspace{-.4cm} \par\nobreak {\small %
    \begin{align*}
      &
        \underbrace{\Pi_\Code{noI}^*\msgAGray{pushI}{m\;y}\msgAGray{popI}{m\;y}\Pi_\Code{noI}^*}_{e_3}\;\;
        \underbrace{\msgA{push}{x}\msgAGray{pushI}{n\;x}}_{e_1}\;\;
        \underbrace{\Pi_\Code{noI}^*}_\text{recursion}\;\;
        \underbrace{\msgA{pop}{x}\msgAGray{popI}{n\;x}}_{e_2}\;\;
        \underbrace{\Pi_\Code{noI}^*}_{e_4}
    \end{align*}}\noindent
  Here, $\pi_\Code{stack}$ is divided into five parts, corresponding
  to the four placeholders and the self-recursive call (i.e., $f\;()$)
  in the template. The synthesizer then re-refines the abstract trace for $i$-fold
  recursion unrolling after erasing all qualifiers from the abstract
  trace: %
  \vspace{-.4cm} \par\nobreak {\small %
    \begin{align*}
      \Pi_\Code{anyA}^*\msgCGray{pushI}\msgCGray{popI}&\;\underline{\epsilon}\;\Pi_\Code{anyA}^* \tag{0-fold unrolling}
      \\\Pi_\Code{anyA}^*\msgCGray{pushI}\msgCGray{popI}\Pi_\Code{anyA}^*\;\underline{\msgCGray{pushI}\msgC{push}}&\underline{\msgCGray{popI}\msgC{pop}}\;\Pi_\Code{anyA}^* \tag{1-fold unrolling}
      \\\Pi_\Code{anyA}^*\msgCGray{pushI}\msgCGray{popI}\Pi_\Code{anyA}^*\;\underline{\msgCGray{pushI}\msgC{push}\msgCGray{pushI}\msgC{push}}&\underline{\msgCGray{popI}\msgC{pop}\msgCGray{popI}\msgC{pop}}\;\Pi_\Code{anyA}^* \tag{2-fold unrolling}
    \end{align*}}\noindent
  This refinement ensures that multiple unfoldings remain well-typed,
  which disallows an incorrect recursion template instantiation, as
  shown below: %
  \vspace{-.4cm} \par\nobreak {\small %
    \begin{align*}
      & \underbrace{\Pi_\Code{noI}^*\msgAGray{pushI}{m\;y}\msgAGray{popI}{m\;y}\Pi_\Code{noI}^*\msgA{push}{x}\msgAGray{pushI}{n\;x}}_{e_3}\;\;
        \underbrace{\Pi_\Code{noI}^*}_{e_1}\;\;
        \underbrace{\Pi_\Code{noI}^*}_\text{recursion}\;\;
        \underbrace{\msgA{pop}{x}\msgAGray{popI}{n\;x}}_{e_2}\;\;
        \underbrace{\Pi_\Code{noI}^*}_{e_4}
    \end{align*}
  } %
  \noindent This instantiation will repeatedly perform $\eff{pop}$
  events without any $\eff{push}$ events; its two-fold recursion
  unrolling will therefore be rejected.  The final recursive program
  is assembled by composing synthesized programs from these traces to
  populate the template placeholders, thereby yielding the following
  program:

  \begin{minted}[fontsize = \footnotesize, escapeinside=??, xleftmargin=10pt, linenos]{ocaml}
let rec f () =
  () ?$\oplus$?
    let (x: int) = int_gen () in push x; (* e1 *)
    f ();
    let (z: int) = pop () in assert (x == z) (*e2 *)
in f ()
  \end{minted}
\end{example}
